\documentclass[10pt,a4paper]{article}
\usepackage[utf8]{inputenc}
\usepackage{amsmath}
\usepackage{amsfonts}
\usepackage{amssymb}
\usepackage{graphicx}
\usepackage{float}
\usepackage{booktabs}
\usepackage{xcolor, soul}
\sethlcolor{green}
\usepackage{enumerate}
\title{GSF-3100 \\
06 Marché obligataire gouvernemental \\
Série d'exercices 5 (Solutions)}

\begin{document}
\maketitle

\begin{itemize}
\item Dans ce document vous avez des questions à choix multiples allant de a) à d).  
\item Il y a seulement un choix possible par question.
\item Ce questionnaire contient des exercices en lien avec le marché obligataire gouvernemental et plus pécisément sur le marché primaire et l'adjudiction d'obligations.
\item Contient 6 questions. 
\end{itemize}
 
\newpage

\section*{Mise en contexte}

Le gouvernement fédéral annonce une nouvelle émission d’obligations standards d’une valeur totale de 20 millions \$.  En réponse, il reçoit les offres de deux distributeurs détaillées ci-bas.  

\begin{table}[H]
\centering
\resizebox{\textwidth}{!}{%
\begin{tabular}{@{}llll@{}}
\toprule
Types d'offres & Offres non concurentielles & Offres concurrentielles #1 & Offres concurrentielles #2 \\ \midrule
Distributeur A & 4 millions \$               & 1 millions \$ à y=2\%        & 1 millions \$ à y=3\%        \\
Distributeur B & 11 millions \$              & 2 millions \$ à y=4\%        & 1 millions \$ à y=5\%        \\ \bottomrule
\end{tabular}%
}
\end{table}

\section*{Exercice 1}
Déterminer en millions de dollars,  le montant total des offres non concurrentielles acceptées.

\begin{enumerate}[a)]
\item 4 millions \$
\item 11 millions \$
\item \hl{15 millions \$}
\item 20 millions \$
\end{enumerate}

\section*{Solution 1}
Nous avons un total de 20 millions \$ (M \$) à répartir entre les offres non concurrentielles et les offres concurrentielles. On débute toujours par attribuer les offres non concurrentielles.

\vspace{1cm}

Le montant total des offres non concurrentielles acceptées:
\begin{itemize}
\item Distributeur A:   4 M \$
\item Distributeur B: 11 M \$
\item Montant total: 15 M \$   
\end{itemize}
Il reste ensuite 5 M \$ à allouer aux offres concurrentielles les plus offrantes (i.e., les offres qui permettent au gouvernement de se financer au plus faible taux possible).

\newpage

\section*{Exercice 2}
Déterminer en millions de dollars,  le montant total des offres concurrentielles acceptées.

\begin{enumerate}[a)]
\item 4 millions \$
\item \hl{5 millions \$}
\item 11 millions \$
\item 15 millions \$
\end{enumerate}

\section*{Solution 2}
Le montant total des offres concurrentielles acceptées:
\begin{itemize}
\item Offre acceptée à y = 2 \%: 1 M \$ (Distributeur A)
\item Offre acceptée à y = 3 \%: 1 M \$ (Distributeur A)
\item Offre acceptée à y = 4\%: 2 M \$ (Distributeur B)
\item Offre acceptée à y = 5 \%: 1 M \$ (Distributeur B)
\item Montant total:                      5 M\$
\end{itemize}

\newpage

\section*{Exercice 3}
Déterminer le taux de rendement des offres non concurrentielles si l'émission est effectuée par le gouvernement canadien.

\begin{enumerate}[a)]
\item 3.00\%
\item 3.50\%
\item \hl{3.60\%}
\item 4.00\%
\end{enumerate}

\section*{Solution 3}
Le taux de rendement des offres non concurrentielles:
\begin{enumerate}
\item Trouver les taux de rendement offert à l'émission aux offres concurrentielles
\begin{itemize}
\item Le distributeur A reçoit 1 M \$ à y = 2 % et 1 M \$ à y = 3 \%. 
\item Le distributeur B reçoit 2 M \$ à y = 4 % et 1 M \$ à y = 5 \%.
\end{itemize}
\item Calculer le taux de rendement moyen pondéré offert à l'émission aux offres non concurrentielles. 
Il faut effectuer le taux moyen pondéré des offres concurrentielles, soit:
\begin{align*}
y=\frac{1}{5} \times 0.02+\frac{1}{5} \times 0.03+\frac{2}{5}\times 0.04 + \frac{1}{5} \times 0.05
\end{align*}
\begin{align*}
y=0.036=3.6\%
\end{align*}
\end{enumerate}


\newpage


\section*{Exercice 4}
Déterminer le taux de coupon des obligations émises si l'émission est effectuée par le gouvernement canadien.

\begin{enumerate}[a)]
\item 3.00\%
\item \hl{3.50\%}
\item 3.60\%
\item 4.00\%
\end{enumerate}

\section*{Solution 4}
Le taux de coupon des obligations émises:
\begin{align*}
TC  = 0,035 = 3,5 \%
\end{align*}
\begin{itemize}
\item Le taux de rendement offert à l'émission aux offres non concurrentielles est de 3,6 \%. Puisque le taux de coupon est arrondi au $\frac{1}{4}$ de pourcentage inférieur,celui-ci sera de 3,5 \%.
\item Explication : 0,6 \% $(\frac{3}{5})$ se situe entre 0,5 \% $(\frac{2}{4})$ et 0,75 \% $(\frac{3}{5})$,  et donc le $\frac{1}{4}$ de pourcentage inférieur se terminera par 0,5 \%. 
\end{itemize}

\newpage

\section*{Exercice 5}
Déterminer le taux de rendement des offres non concurrentielles si l'émission est effectuée par par le gouvernement américain.

\begin{enumerate}[a)]
\item 3.50\%
\item 4.00\%
\item 4.50\%
\item \hl{5.00\%}
\end{enumerate}

\section*{Solution 5}
\begin{itemize}
\item Gouvernement américain: Adjudication à prix unique
\item Les offres non concurrentielles et concurrentielles sont offertes à l'émission à un taux unique de rendement. Il s'agit du dernier taux de rendement accepté qui permet de compléter l'adjudication.
\end{itemize}
Le taux de rendement des offres non concurrentielles:

\vspace{0.5cm}

Le dernier taux de rendement accepté est 5 \% pour un montant de 1 M \$.  Donc, le taux de rendement des offres concurrentielles et non concurrentielles est de y = 5 \%. 


\newpage

\section*{Exercice 6}
Déterminer le taux de coupon des obligations émises si l'émission est effectuée par par le gouvernement américain.

\begin{enumerate}[a)]
\item 4.50\%
\item \hl{4.75\%}
\item 5.00\%
\item 5.25\%
\end{enumerate}

\section*{Solution 6}
Le taux de coupon des obligations émises:
\begin{align*}
TC  = 0,0475 = 4,75 \%
\end{align*}

Le taux de rendement de 5 \% arrondi au quart de pourcentage inférieur donne un taux de coupon pour toutes les obligations de 4,75 \%.
\end{document}
